\uuid{7maZ}
\titre{Notion d'équation différentielle}
\auteur{bodin ; croizat ; sacre}
\organisation{exo7}
\licence{CC BY-NC-SA 4.0}
\datecreate{2022-03-08}
\niveau{L1}
\module{Analyse}
\chapitre{Equation différentielle}
\sousChapitre{Résolution d'équation différentielle du premier ordre}
\difficulte{3}
\theme{Notion d'équation différentielle}
\exo7qcmid{hilisit.equations-differentielles.024}

\contenu{
\begin{question}
\qclassification{225.01}
Soit $f$ une solution de l'équation différentielle $y' = \cos(x) y$, vérifiant $f(\frac\pi3)=3$. On considère la courbe représentative de \(f\).
Quelles sont les affirmations vraies ?
\begin{answers}
  \bad{La tangente en $x=\frac\pi3$ a pour équation $y=\frac32x + 3 $.}
  \good{La tangente en $x=\frac\pi3$ a pour équation $y=\frac32(x-\frac\pi3) + 3$.}
  \good{La tangente en $x=\frac\pi2$ est horizontale.}
  \bad{La tangente en $x=\frac\pi3$ est horizontale.}
\end{answers}
\begin{explanations}
En $x=\frac\pi2$, par l'équation différentielle on a $f'(\frac\pi2) = 0$ (car $\cos\frac\pi2=0$), donc la tangente est horizontale.
En $x=\frac\pi3$, on obtient $f'(\frac\pi3) = \cos(\frac\pi3) y(\frac\pi3) = \frac12 \times 3 = \frac32$, donc la pente de la tangente en $x=\frac\pi3$ est $\frac32$. Cette tangente passe par le point $(\frac\pi3,3)$ donc son équation est $y=\frac32(x-\frac\pi3) + 3$.
\end{explanations}
\end{question}
}
